% Options for packages loaded elsewhere
\PassOptionsToPackage{unicode}{hyperref}
\PassOptionsToPackage{hyphens}{url}
\documentclass[
]{article}
\usepackage{xcolor}
\usepackage[margin=1in]{geometry}
\usepackage{amsmath,amssymb}
\setcounter{secnumdepth}{-\maxdimen} % remove section numbering
\usepackage{iftex}
\ifPDFTeX
  \usepackage[T1]{fontenc}
  \usepackage[utf8]{inputenc}
  \usepackage{textcomp} % provide euro and other symbols
\else % if luatex or xetex
  \usepackage{unicode-math} % this also loads fontspec
  \defaultfontfeatures{Scale=MatchLowercase}
  \defaultfontfeatures[\rmfamily]{Ligatures=TeX,Scale=1}
\fi
\usepackage{lmodern}
\ifPDFTeX\else
  % xetex/luatex font selection
\fi
% Use upquote if available, for straight quotes in verbatim environments
\IfFileExists{upquote.sty}{\usepackage{upquote}}{}
\IfFileExists{microtype.sty}{% use microtype if available
  \usepackage[]{microtype}
  \UseMicrotypeSet[protrusion]{basicmath} % disable protrusion for tt fonts
}{}
\makeatletter
\@ifundefined{KOMAClassName}{% if non-KOMA class
  \IfFileExists{parskip.sty}{%
    \usepackage{parskip}
  }{% else
    \setlength{\parindent}{0pt}
    \setlength{\parskip}{6pt plus 2pt minus 1pt}}
}{% if KOMA class
  \KOMAoptions{parskip=half}}
\makeatother
\usepackage{color}
\usepackage{fancyvrb}
\newcommand{\VerbBar}{|}
\newcommand{\VERB}{\Verb[commandchars=\\\{\}]}
\DefineVerbatimEnvironment{Highlighting}{Verbatim}{commandchars=\\\{\}}
% Add ',fontsize=\small' for more characters per line
\usepackage{framed}
\definecolor{shadecolor}{RGB}{248,248,248}
\newenvironment{Shaded}{\begin{snugshade}}{\end{snugshade}}
\newcommand{\AlertTok}[1]{\textcolor[rgb]{0.94,0.16,0.16}{#1}}
\newcommand{\AnnotationTok}[1]{\textcolor[rgb]{0.56,0.35,0.01}{\textbf{\textit{#1}}}}
\newcommand{\AttributeTok}[1]{\textcolor[rgb]{0.13,0.29,0.53}{#1}}
\newcommand{\BaseNTok}[1]{\textcolor[rgb]{0.00,0.00,0.81}{#1}}
\newcommand{\BuiltInTok}[1]{#1}
\newcommand{\CharTok}[1]{\textcolor[rgb]{0.31,0.60,0.02}{#1}}
\newcommand{\CommentTok}[1]{\textcolor[rgb]{0.56,0.35,0.01}{\textit{#1}}}
\newcommand{\CommentVarTok}[1]{\textcolor[rgb]{0.56,0.35,0.01}{\textbf{\textit{#1}}}}
\newcommand{\ConstantTok}[1]{\textcolor[rgb]{0.56,0.35,0.01}{#1}}
\newcommand{\ControlFlowTok}[1]{\textcolor[rgb]{0.13,0.29,0.53}{\textbf{#1}}}
\newcommand{\DataTypeTok}[1]{\textcolor[rgb]{0.13,0.29,0.53}{#1}}
\newcommand{\DecValTok}[1]{\textcolor[rgb]{0.00,0.00,0.81}{#1}}
\newcommand{\DocumentationTok}[1]{\textcolor[rgb]{0.56,0.35,0.01}{\textbf{\textit{#1}}}}
\newcommand{\ErrorTok}[1]{\textcolor[rgb]{0.64,0.00,0.00}{\textbf{#1}}}
\newcommand{\ExtensionTok}[1]{#1}
\newcommand{\FloatTok}[1]{\textcolor[rgb]{0.00,0.00,0.81}{#1}}
\newcommand{\FunctionTok}[1]{\textcolor[rgb]{0.13,0.29,0.53}{\textbf{#1}}}
\newcommand{\ImportTok}[1]{#1}
\newcommand{\InformationTok}[1]{\textcolor[rgb]{0.56,0.35,0.01}{\textbf{\textit{#1}}}}
\newcommand{\KeywordTok}[1]{\textcolor[rgb]{0.13,0.29,0.53}{\textbf{#1}}}
\newcommand{\NormalTok}[1]{#1}
\newcommand{\OperatorTok}[1]{\textcolor[rgb]{0.81,0.36,0.00}{\textbf{#1}}}
\newcommand{\OtherTok}[1]{\textcolor[rgb]{0.56,0.35,0.01}{#1}}
\newcommand{\PreprocessorTok}[1]{\textcolor[rgb]{0.56,0.35,0.01}{\textit{#1}}}
\newcommand{\RegionMarkerTok}[1]{#1}
\newcommand{\SpecialCharTok}[1]{\textcolor[rgb]{0.81,0.36,0.00}{\textbf{#1}}}
\newcommand{\SpecialStringTok}[1]{\textcolor[rgb]{0.31,0.60,0.02}{#1}}
\newcommand{\StringTok}[1]{\textcolor[rgb]{0.31,0.60,0.02}{#1}}
\newcommand{\VariableTok}[1]{\textcolor[rgb]{0.00,0.00,0.00}{#1}}
\newcommand{\VerbatimStringTok}[1]{\textcolor[rgb]{0.31,0.60,0.02}{#1}}
\newcommand{\WarningTok}[1]{\textcolor[rgb]{0.56,0.35,0.01}{\textbf{\textit{#1}}}}
\usepackage{graphicx}
\makeatletter
\newsavebox\pandoc@box
\newcommand*\pandocbounded[1]{% scales image to fit in text height/width
  \sbox\pandoc@box{#1}%
  \Gscale@div\@tempa{\textheight}{\dimexpr\ht\pandoc@box+\dp\pandoc@box\relax}%
  \Gscale@div\@tempb{\linewidth}{\wd\pandoc@box}%
  \ifdim\@tempb\p@<\@tempa\p@\let\@tempa\@tempb\fi% select the smaller of both
  \ifdim\@tempa\p@<\p@\scalebox{\@tempa}{\usebox\pandoc@box}%
  \else\usebox{\pandoc@box}%
  \fi%
}
% Set default figure placement to htbp
\def\fps@figure{htbp}
\makeatother
\usepackage{svg}
\setlength{\emergencystretch}{3em} % prevent overfull lines
\providecommand{\tightlist}{%
  \setlength{\itemsep}{0pt}\setlength{\parskip}{0pt}}
\usepackage{bookmark}
\IfFileExists{xurl.sty}{\usepackage{xurl}}{} % add URL line breaks if available
\urlstyle{same}
\hypersetup{
  hidelinks,
  pdfcreator={LaTeX via pandoc}}

\author{}
\date{\vspace{-2.5em}}

\begin{document}

\href{https://CRAN.R-project.org/package=ModTools}{\pandocbounded{\includegraphics[keepaspectratio,alt={CRAN status}]{https://www.r-pkg.org/badges/version-last-release/ModTools}}}
\href{https://CRAN.R-project.org/package=ModTools}{\pandocbounded{\includegraphics[keepaspectratio,alt={downloads}]{https://cranlogs.r-pkg.org/badges/grand-total/ModTools}}}
\href{https://CRAN.R-project.org/package=ModTools}{\pandocbounded{\includegraphics[keepaspectratio,alt={downloads}]{http://cranlogs.r-pkg.org/badges/last-week/ModTools}}}
\href{https://www.gnu.org/licenses/old-licenses/gpl-2.0.en.html}{\pandocbounded{\includesvg[keepaspectratio]{https://img.shields.io/badge/License-GPL\%20v2+-blue.svg}}}
\href{https://lifecycle.r-lib.org/articles/stages.html}{\pandocbounded{\includesvg[keepaspectratio]{https://img.shields.io/badge/lifecycle-maturing-blue.svg}}}
\href{https://github.com/AndriSignorell/ModTools/actions}{\pandocbounded{\includesvg[keepaspectratio]{https://github.com/AndriSignorell/ModTools/workflows/R-CMD-check/badge.svg}}}
\href{https://andrisignorell.github.io/ModTools/}{\pandocbounded{\includesvg[keepaspectratio]{https://github.com/AndriSignorell/ModTools/workflows/pkgdown/badge.svg}}}

\section{Tools for Building Regression and Classification
Models}\label{tools-for-building-regression-and-classification-models}

A unified interface to the most common regression and classification
algorithms --- including random forests, neural networks, decision
trees, and support vector machines --- blending them into a single,
consistent workflow. Complements model fitting with tools for
validation, tuning, variable importance, and ROC analysis.

There is a rich selection of R packages implementing algorithms for
classification and regression tasks out there. The authors legitimately
take the liberty to tailor the function interfaces according to their
own taste and needs. For us other users, however, this often results in
struggling with user interfaces, some of which are rather weird - to put
it mildly - and almost always different in terms of arguments and result
structures. ModTools pursues the goal of offering uniform handling for
the most important regression and classification models in applied data
analyses. The function FitMod() is designed as a simple and consistent
interface to these original functions while maintaining the flexibility
to pass on all possible arguments. print, plot, summary and predict
operations can so be carried out following the same logic. The results
will again be reshaped to a reasonable standard. For all the functions
of this package Google styleguides are used as naming rules (in absence
of convincing alternatives). The 'BigCamelCase' style has been
consequently applied to functions borrowed from contributed R packages
as well.

As always: Feedback, feature requests, bug reports and other suggestions
are welcome! Please report problems to
\href{https://github.com/AndriSignorell/ModTools/issues}{GitHub issues
tracker} (preferred) or directly to the maintainer.

\subsection{Installation}\label{installation}

You can install the released version of \textbf{ModTools} from
\href{https://CRAN.R-project.org}{CRAN} with:

\begin{Shaded}
\begin{Highlighting}[]
\FunctionTok{install.packages}\NormalTok{(}\StringTok{"ModTools"}\NormalTok{)}
\end{Highlighting}
\end{Shaded}

And the development version from GitHub with:

\begin{Shaded}
\begin{Highlighting}[]
\ControlFlowTok{if}\NormalTok{ (}\SpecialCharTok{!}\FunctionTok{require}\NormalTok{(}\StringTok{"remotes"}\NormalTok{)) }\FunctionTok{install.packages}\NormalTok{(}\StringTok{"remotes"}\NormalTok{)}
\NormalTok{remotes}\SpecialCharTok{::}\FunctionTok{install\_github}\NormalTok{(}\StringTok{"AndriSignorell/ModTools"}\NormalTok{)}
\end{Highlighting}
\end{Shaded}

\section{Warning}\label{warning}

\textbf{Warning:} This package is still under development. Although the
code seems meanwhile quite stable, until release of version 1.0 you
should be aware that everything in the package might be subject to
change. Backward compatibility is not yet guaranteed. Functions may be
deleted or renamed and new syntax may be inconsistent with earlier
versions. By release of version 1.0 the ``deprecated-defunct process''
will be installed.

\section{Authors}\label{authors}

Andri Signorell\\
Helsana Versicherungen AG, Health Sciences, Zurich\\
HWZ University of Applied Sciences in Business Administration Zurich.

R is a community project. This can be seen from the fact that this
package includes R source code and/or documentation previously published
by \href{https://andrisignorell.github.io/ModTools/authors.html}{various
authors and contributors}. Special thanks go to Beat Bruengger, Mathias
Frueh, Daniel Wollschlaeger for their valuable contributions and
testing. The good things come from all these guys, any problems are
likely due to my tweaking. Thank you all!

\textbf{Maintainer:} Andri Signorell

\section{Examples}\label{examples}

\begin{Shaded}
\begin{Highlighting}[]
\FunctionTok{library}\NormalTok{(ModTools)}
\end{Highlighting}
\end{Shaded}

\begin{Shaded}
\begin{Highlighting}[]
\FunctionTok{demo}\NormalTok{(bin\_mod, }\AttributeTok{package =} \StringTok{"ModTools"}\NormalTok{)}
\end{Highlighting}
\end{Shaded}

\begin{Shaded}
\begin{Highlighting}[]
\FunctionTok{demo}\NormalTok{(mult\_mod, }\AttributeTok{package =} \StringTok{"ModTools"}\NormalTok{)}
\end{Highlighting}
\end{Shaded}

\begin{Shaded}
\begin{Highlighting}[]
\FunctionTok{demo}\NormalTok{(num\_mod, }\AttributeTok{package =} \StringTok{"ModTools"}\NormalTok{)}
\end{Highlighting}
\end{Shaded}


\end{document}
